% !TEX TS-program = xelatex
% !TEX encoding = UTF-8 Unicode
% !Mode:: "TeX:UTF-8"

\documentclass{resume}
\usepackage{zh_CN-Adobefonts_external} % Simplified Chinese Support using external fonts (./fonts/zh_CN-Adobe/)
%\usepackage{zh_CN-Adobefonts_internal} % Simplified Chinese Support using system fonts
\usepackage{linespacing_fix} % disable extra space before next section
\usepackage{cite}
\usepackage[hidelinks]{hyperref} % 可点击链接（无边框）

\begin{document}
\pagenumbering{gobble} % suppress displaying page number

\name{张震}
% {E-mail}{mobilephone}{homepage}
% be careful of _ in emaill address
\contactInfo{(+86) 17600804803}{zhangz1728@126.com}{\href{https://markz.fun}{markz.fun}}{}
% {E-mail}{mobilephone}
% keep the last empty braces!
%\contactInfo{xxx@yuanbin.me}{(+86) 131-221-87xxx}{}

% 照片（右上角）
\photo{./photo/photo.jpg}

\section{教育背景}
\datedsubsection{\textbf{兰州理工大学},计算机技术,\textit{硕士}}{2022.9 - 2025.6}
全日制, 计算机与通信学院, CAA Student member

\section{工作经历}
\datedsubsection{\textbf{百度在线网络技术（北京）有限公司}}{2025.7 - 至今}
\role{客户端开发工程师}{}
负责落地页、激励视频与 AI 工程化方向开发。

\datedsubsection{\textbf{落地页}}{}
\begin{onehalfspacing}
  \begin{itemize}
    \item \textbf{广告落地页容器}：深入理解落地页容器框架的调起、宿主与扩展机制，主导多套容器统一重构，统一页面调起与容器能力，新业务接入复用同一容器。
    \item \textbf{预渲染}：针对落地页首屏等待问题，设计点击前预加载、点击后复用实例的预渲染链路，未命中自动回退普通加载；半屏场景预渲染生效比例由 66\% 提升至 99.6\%。
    \item \textbf{半屏落地页框架}：基于统一容器落地半屏形态，复用主容器与预渲染能力，支持半屏内跳转、下载、拖拽关闭等交互（获 2026Q1 百度小赞·重大项目）。
  \end{itemize}
\end{onehalfspacing}

\datedsubsection{\textbf{激励视频}}{}
\begin{onehalfspacing}
  \begin{itemize}
    \item \textbf{激励视频组件化重构}：针对单页面逻辑集中的问题，将播放、倒计时、任务、弹窗、UI 等按职责拆分为可复用组件，页面只做编排、组件间通过统一接口通信，降低耦合，提升迭代与复用效率。
    \item \textbf{弹窗体系改造}：统一激励视频弹窗管理，覆盖领奖、挽留、转化引导等多类弹窗，支持自定义布局样式；处理弹窗与播放、倒计时的联动及展示优先级、频控，保障完播与领奖链路，并解决相关内存泄漏问题。
  \end{itemize}
\end{onehalfspacing}

\datedsubsection{\textbf{AI 工程化改造}}{}
\begin{onehalfspacing}
  \begin{itemize}
    \item \textbf{团队 Agent 插件体系}：跟进 Agent 生态前沿，从 0 搭建团队 Agent 插件仓库框架（rules、skills、plugins、memory 四层资产）。提出插件化维护方式，并适配 Claude Code、Codex、厂内 agent 开发工具等多个宿主，支撑资产版本化分发、安装与回滚。
    \item \textbf{Harness 资产指标}：建设团队 Harness 资产指标插件，多 runtime hooks 采集与事件归一化，按 Skill 聚合调用次数、耗时、成功率并日报上报，为 AI 工具效果评估与资产迭代提供数据支撑。
    \item \textbf{团队 Agent Workspace（提出、主导，推进中）}：面向项目组织 Agent 研发所需的长期操作空间——五层资产（代码、知识、环境、工具、流程）+ 三平面（上下文、能力、治理）+ 四项治理边界（Repo 路由、事实 Owner、权限边界、发布节奏）；稳定资产经 Owner 评审发布，任务现场只回写经 Review 的增量。
  \end{itemize}
\end{onehalfspacing}

\datedsubsection{\textbf{通用组件建设}}{}
\begin{onehalfspacing}
  \begin{itemize}
    \item \textbf{通用按钮统一}：作为负责人统一多业务按钮实现与接入规范，推动通用按钮在多场景扩覆盖、统一按钮出口；推动通用按钮模型下沉 lib\_core，降低多业务重复开发成本；主导实验与 LR 协同，实验验证大盘持平、CVR +0.10\%。
    \item \textbf{多宿主依赖隔离}：针对不同宿主对 PAG 等三方能力支持不一致的问题，设计依赖隔离与自动降级方案（宿主具备 PAG 时走实现、缺失时自动降级），避免构建与运行风险。
    \item \textbf{组件与弹窗优先级管理（SDK 化）}：抽象统一的优先级模型与互斥裁决机制，组件与弹窗注册时声明优先级与互斥关系，由管理器统一裁决展示、关闭与降级顺序；设计通用接口与接入规范，把各业务重复的互斥逻辑收敛为 SDK 公共能力，并与激励视频弹窗优先级链、组件互斥治理统一为一套模型。
  \end{itemize}
\end{onehalfspacing}

\datedsubsection{\textbf{私信商业化（2025Q4，获季度 Award、百度小赞）}}{}
\begin{onehalfspacing}
  \begin{itemize}
    \item \textbf{插件化接入与分享}：组内首个承接私信项目，快速吃透私信插件化接入原理，完成私信作为广告透传载体的 UI 侧 0 到 1 开发与上线；沉淀接入方法并在组内分享，帮助同事快速接入开发。
    \item \textbf{跨团队协作}：打通与 IM 团队合作链路，推动协作模式从“对方出方案 + 我方开发”演进为“我方出方案 + 对方 review”；协调 7-8 个模块联调，推动消息收发 AI 化。
  \end{itemize}
\end{onehalfspacing}

\section{个人项目}
\datedsubsection{\textbf{JSONUtils 专业版 —— 全 AI 开发 JSON 工具}}{}
\textbf{GitHub}：\underline{\href{https://github.com/zhang99667/jsonUtil}{github.com/zhang99667/jsonUtil}} \quad \textbf{在线}：\underline{\href{https://jsonutils.markz.fun}{jsonutils.markz.fun}}
  \begin{itemize}
    \item \textbf{AI 驱动的大型项目研发}：独立开发 JSON 处理与接口排查工具 JSONUtils（TypeScript / React + Spring Boot + Electron），12w+ 行代码、1400+ 次提交；需求拆解、方案、编码、自测到 CR 全流程 AI 参与。
    \item \textbf{多 Agent 研发体系}：搭建 Codex / Claude / Copilot 三端 Agent 协作体系——AGENTS.md / CLAUDE.md 双入口、rules 与 skills、explorer / worker / verifier / ai-infra-auditor 子 Agent 委派与治理 MCP。
    \item \textbf{AI 治理与评测闭环}：搭建 AI 资产治理体系——38 条 eval case、grader 校准与 append-only ledger，区分行为覆盖与组件证据；配套 6 份 playbook 与十余项 CI 门禁。
  \end{itemize}


\section{实习经历}
\datedsubsection{\textbf{百度在线网络技术（北京）有限公司}}{2024.12 - 2025.4}
\role{客户端开发工程师}{}
\begin{onehalfspacing}
  \begin{itemize}
    \item \textbf{WebView 预加载与预渲染优化}：设计 WebView 复用池与生命周期模型，结合空闲时机在后台创建 WebView，初始化耗时由约 300ms 降至约 70ms（-77\%），提升广告页面加载速度与打开体验。
    \item \textbf{Android 组件管理优化}：定义统一组件规范接口并引入组件管理器，收敛组件互斥与交互逻辑；将单例持有的 HashMap 改为随界面生命周期管理，消除潜在内存泄漏。
  \end{itemize}
\end{onehalfspacing}

\section{科研经历}
\begin{itemize}[parsep=0.2ex]
  \item The Influence of Opinion Leaders on Public Opinion Spread and Control Strategies in Online Social Networks. (SCI 二区，DOI: 10.1109/TCSS.2025.3537701)
  \item Information propagation under the coupling of behavioral and influential propagation (SCI 三区，审稿中)
  \item 基于大数据的“人肉搜索”建模、分析与控制，国家自然科学地区基金(No.61863025)
  \item 网络信息相变及其传播动力学特征分析研究，国家自然科学地区基金(No.62266030)
\end{itemize}

\section{技术能力}
% increase linespacing [parsep=0.5ex]
\begin{itemize}[parsep=0.2ex]
  \item 具备扎实的\textbf{Java基础}, 熟悉常用集合\textbf{ArrayList、HashMap}等底层原理, 熟悉基本的\textbf{数据结构与算法}
  \item 熟悉JUC掌握volatile、synchronized的使用以及CAS、线程池等相关知识
  \item 熟悉JVM运行时内存区域、引用类型、垃圾回收算法、类加载过程等
  \item 熟练使用\textbf{Spring、SpringMVC、MyBatis与MyBatisPlus}等主流开发框架以及\textbf{使用SpringBoot整合第三方框架}, 了解\textbf{IOC、AOP原理}
  \item 熟练使用\textbf{MySQL关系型数据库}, 了解MySQL常见的知识, 如锁、事务、索引等。熟悉SQL语句编写, 了解常见的SQL优化
  \item 熟练使用\textbf{非关系型数据库Redis}, 了解\textbf{缓存雪崩、缓存穿透、缓存击穿}、持久化方式等
  \item 熟练使用Maven、Git等工具。熟练使用Swagger、Apifox、Postman等工具进行接口测试, 熟悉Restful风格接口开发
  \item 熟悉\textbf{Linux常用指令}, 可以在Linux上\textbf{搭建开发环境和项目部署}
  \item \textbf{常用编程语言}: Java~\,/\,~Kotlin~\,/\,~Python
  \item \textbf{AI 工程化}：Agent Plugin~\,/\,~Skill~\,/\,~Hook 机制与多宿主适配，Skill 埋点与效果评估；熟悉 Claude Code、Codex 等 Agent 运行时
\end{itemize}

\section{工作荣誉}
\noindent 2026Q3 百度小赞 · \textbf{AI 工程化先锋}；2026Q1 百度小赞 · \textbf{落地页半屏重大项目}；2025Q4 季度 Award + 百度小赞 · \textbf{私信商业化（广告透传）}。

%% Reference
%\newpage
%\bibliographystyle{IEEETran}
%\bibliography{mycite}
\end{document}
